% =====================================================================
%  1.7 Implementação dos sistemas de abastecimento de água
%  (Atividades 3.2.2.3 e 3.2.2.6)
% =====================================================================
\section{Implementação dos sistemas de abastecimento de água}
\label{sec:abastecimento}

O IAT implantará ou reabilitará 83 sistemas coletivos com captação por poço
tubular, atendendo 5.810 domicílios, entre 2028 e 2033. A
\autoref{fig:abastecimento} mostra o percurso completo, da preparação à
entrega do sistema à entidade gestora.

\begin{figure}[htbp]
  \centering
  \caption{Implementação do abastecimento de água, do começo ao fim}
  \label{fig:abastecimento}
  \fluxograma{fluxograma-abastecimento}
  \fonte{Elaborado pelo IAT (2026) com base em \citeonline{mop2026} e no POA 2026--2027.}
\end{figure}

% ---------------------------------------------------------------------
\subsection{Captação, projeto e licenciamento}
\label{subsec:captacao}

A captação começa pela verificação da titularidade das áreas e pelo estudo
hidrogeológico, que define o ponto do poço. Depois da perfuração, o teste de
vazão e a análise da água orientam o projeto de tratamento, reservação e
distribuição, que só vai à obra com outorga e licenças
(\autoref{fig:captacao-projeto}; \autoref{qua:etapas-3223}).

\begin{figure}[htbp]
  \centering
  \caption{Captação subterrânea, projeto e licenciamento}
  \label{fig:captacao-projeto}
  \fluxograma{fluxograma-captacao-projeto}
  \fonte{Elaborado pelo IAT (2026) com base em \citeonline{mop2026}.}
\end{figure}

\pendencia{O MOP prevê ``contratação integrada'', mas a define como
semi-integrada (Lei nº~14.133/2021, art.~6º, XXXIII). O POA prevê aquisição
de materiais pelo IAT e contratação separada dos projetos. Definir o modelo
de execução e quem executa as obras com os materiais adquiridos.}

% ---------------------------------------------------------------------
\subsection{Execução das obras}
\label{subsec:obras}

A obra compreende casa química, reservatório, parte eletromecânica, rede de
distribuição e hidrômetros. Uma empresa supervisora acompanha a execução, e a
fiscalização do IAT verifica a conformidade com o projeto antes de emitir o
atesto de recebimento (\autoref{fig:obras}).

\begin{figure}[htbp]
  \centering
  \caption{Execução, supervisão e fiscalização das obras, por responsável}
  \label{fig:obras}
  \fluxograma{fluxograma-obras}
  \fonte{Elaborado pelo IAT (2026) com base em \citeonline{mop2026}.}
\end{figure}

\pendencia{Os lotes de materiais do POA (25 + 30 = 55 unidades) não cobrem os
83 sistemas. A quantidade de hidrômetros ainda está como ``XXX'' no MOP.}

% ---------------------------------------------------------------------
\subsection{Entrega e operação}
\label{subsec:entrega}

Concluída a obra, o sistema passa por pré-operação e testes, os operadores
locais são capacitados e a entidade gestora assina o aceite. O IAT mantém o
suporte técnico e avalia o desempenho a cada semestre até o Termo de
Recebimento Definitivo (\autoref{fig:entrega}).

\begin{figure}[htbp]
  \centering
  \caption{Transferência, operação e encerramento dos sistemas de abastecimento}
  \label{fig:entrega}
  \fluxograma{fluxograma-entrega}
  \fonte{Elaborado pelo IAT (2026) com base em \citeonline{mop2026}.}
\end{figure}

% ---------------------------------------------------------------------
\subsection{Programa Sanepar Rural}
\label{subsec:sanepar-rural}

Nos municípios com contrato de concessão com a Sanepar, o Programa Sanepar
Rural implantará 222 sistemas coletivos entre 2027 e 2031, com R\$ 67,1
milhões de contrapartida \cite{custo2026}. A Sanepar faz os estudos e
projetos e fornece os materiais. O município executa as obras, e a comunidade
forma sua associação (\autoref{fig:sanepar-rural}).

\begin{figure}[htbp]
  \centering
  \caption{Programa Sanepar Rural, por responsável}
  \label{fig:sanepar-rural}
  \fluxograma{fluxograma-sanepar-rural}
  \fonte{Elaborado pelo IAT (2026) com base em \citeonline{mop2026}.}
\end{figure}

\pendencia{O Quadro 31 do MOP atribui ao Sanepar Rural a meta total do
Componente~3 (59.780 pessoas), que inclui os sistemas do IAT. A parcela da
Sanepar é de cerca de 43.500 pessoas. Citar o manual do programa com título,
versão e data.}
